\section{Methodology}
\label{sec:method}

In this section, we present the mathematical formulation of \textbf{Unitary Geodesic Routing (UGR)}, a novel test-time model ensembling framework. UGR models the state of the ensembling weights entirely on a curved hyperspherical manifold, performing state updates via closed-form geodesic rotations and projecting onto the probability simplex natively.

\subsection{Problem Setup}
Let $X_t$ be a sequential stream of input queries at time step $t \in \{1, 2, \dots\}$. We consider a pre-trained deep neural network backbone consisting of $L$ layers. The first $L_{\text{frozen}}$ layers are frozen and shared, while the remaining layers $l \in \{L_{\text{frozen}}+1, \dots, L\}$ are adapted using $K$ parallel task-specific expert adapters. 

For each layer $l$, let $h_t^{(l-1)} \in \mathbb{R}^D$ be the intermediate representation entering the adapted layer. Our goal is to dynamically compute a set of ensembling weights $\boldsymbol{\alpha}_t^{(l)} \in \Delta^{K-1}$ to blend the expert activations such that the output is:
\begin{equation}
h_t^{(l)} = h_t^{(l-1)} W_{\text{base}}^{(l)} + \sum_{k=1}^K \alpha_{k, t}^{(l)} \left( h_t^{(l-1)} A_k^{(l)} B_k^{(l)} \right)
\label{eq:forward}
\end{equation}
where $W_{\text{base}}^{(l)}$ represents the base layer parameters, and $A_k^{(l)}, B_k^{(l)}$ are the down- and up-projection weights of the $k$-th expert adapter. The weights $\boldsymbol{\alpha}_t^{(l)}$ must satisfy the probability simplex constraints:
\begin{equation}
\Delta^{K-1} = \left\{ \boldsymbol{\alpha} \in \mathbb{R}^K : \sum_{k=1}^K \alpha_k = 1 \quad \text{and} \quad \alpha_k \ge 0 \quad \forall k \right\}
\end{equation}

\subsection{Spherical State Representation and Fisher-Rao Born Mapping}
Rather than maintaining an unconstrained state $\mathbf{z} \in \mathbb{R}^K$ and applying a Softmax projection, we model the ensembling state at time step $t$ and adapted layer $l$ as a unit vector on the $(K-1)$-dimensional hypersphere:
\begin{equation}
\mathbf{s}_t^{(l)} \in \mathbb{S}^{K-1} = \left\{ \mathbf{s} \in \mathbb{R}^K : \|\mathbf{s}\|_2 = 1 \right\}
\end{equation}

To project this spherical state onto the probability simplex $\Delta^{K-1}$ in a Softmax-free, scale-preserving manner, we define the ensembling weights $\alpha_{k, t}^{(l)}$ coordinate-wise as the squared magnitudes of the state vector:
\begin{equation}
\alpha_{k, t}^{(l)} = \left(s_{k, t}^{(l)}\right)^2
\label{eq:born}
\end{equation}

This formulation is primarily grounded in the foundational principles of \textbf{Information Geometry}, while sharing an elegant mathematical analogy with Born's rule of quantum mechanics (where probability density is given by the squared magnitude of the wave function). We explicitly note that the term ``Unitary'' in Unitary Geodesic Routing refers strictly to the unit-norm constraint of our ensembling state vector residing on the unit hypersphere ($\|\mathbf{s}\|_2 = 1$), rather than any quantum unitary operators. In information geometry, the square-root map $s_k = \sqrt{\alpha_k}$ is the standard homeomorphism (historically termed the Hellinger or Bhattacharyya mapping) mapping the probability simplex $\Delta^{K-1}$ to the positive orthant of the unit hypersphere $\mathbb{S}^{K-1}_+$. Under this mapping, the celebrated \textbf{Fisher-Rao Riemannian metric} on the probability simplex corresponds exactly to the standard round Riemannian metric (geodesic arc-length) on the sphere. Thus, by restricting updates to geodesics on the hypersphere and mapping to the simplex via Equation \eqref{eq:born}, UGR's trajectories represent exact, closed-form Fisher-Rao geodesic flows on the probability simplex. We verify that Equation \eqref{eq:born} guarantees the simplex constraints natively:
\begin{equation}
\sum_k \alpha_{k, t}^{(l)} = \sum_k \left(s_{k, t}^{(l)}\right)^2 = \left\|\mathbf{s}_t^{(l)}\right\|_2^2 = 1, \quad \alpha_{k, t}^{(l)} \ge 0 \quad \forall k
\end{equation}
By bypassing the Softmax layer, we prevent representational compression at the simplex boundaries, preserving the scale and norm of intermediate activations.

\textbf{Sign Symmetry and Positive Orthant Persistence.} Equation \eqref{eq:born} implies that the ensembling weights $\boldsymbol{\alpha}$ are invariant to coordinate-wise sign flips of the state $\mathbf{s}$. However, in our framework, the bottom-up target vector $\mathbf{w}_t^{(l)}$ is constructed from a non-negative probability distribution normalized to the sphere, meaning $\mathbf{w}_t^{(l)}$ resides strictly in the positive orthant of the hypersphere:
\begin{equation}
\mathbb{S}^{K-1}_+ = \left\{ \mathbf{s} \in \mathbb{S}^{K-1} \mid s_k \ge 0 \quad \forall k \right\}
\end{equation}
Since we initialize the state uniformly as $\mathbf{s}_1 = \frac{1}{\sqrt{K}}\mathbf{1} > 0$ and perform geodesic rotations along the shortest great-circle path (which connects points in the positive orthant), the entire ensembling trajectory $\mathbf{s}_t^{(l)}$ is mathematically guaranteed to remain strictly within the positive orthant $\mathbb{S}^{K-1}_+$. Consequently, coordinate signs never flip, and antipodal path ambiguities do not arise during geodesic updates.

\subsection{Bottom-Up Target Vector Construction}
At each adapted layer $l$, the incoming activation $h_t^{(l-1)}$ is evaluated against pre-computed task centroids $\mu_k^{(l)} \in \mathbb{R}^D$ ($k \in \{1, \dots, K\}$) to extract a localized, bottom-up task-similarity signal. We compute the cosine similarity and apply a Softmax with a tuning temperature $\tau$:
\begin{equation}
e_{k, t}^{(l)} = \frac{\exp\left( \text{sim}(h_t^{(l-1)}, \mu_k^{(l)}) / \tau \right)}{\sum_{j=1}^K \exp\left( \text{sim}(h_t^{(l-1)}, \mu_j^{(l)}) / \tau \right)}
\label{eq:target_softmax}
\end{equation}
where $\mathbf{e}_t^{(l)} \in [0, 1]^K$ represents the stateless target ensembling distribution. Note that we employ a localized Softmax over similarity scores strictly to extract the target signal $\mathbf{e}_t^{(l)}$, which aligns our target space with baseline environments for a fair comparison. Crucially, our \emph{state representation}, \emph{recurrent trajectory updates}, and \emph{simplex projections} are entirely Softmax-free, avoiding the representational scale distortions and boundary clipping of baseline methods. Furthermore, UGR is fully compatible with completely Softmax-free target constructions: for instance, one can replace the target Softmax in Equation \eqref{eq:target_softmax} with a simple rectified linear unit (ReLU) and $L_1$-normalization:
\begin{equation}
e_{k, t}^{(l)} = \frac{\text{ReLU}\left( \text{sim}(h_t^{(l-1)}, \mu_k^{(l)}) \right)}{\sum_{j=1}^K \text{ReLU}\left( \text{sim}(h_t^{(l-1)}, \mu_j^{(l)}) \right)}
\label{eq:target_relu}
\end{equation}
Because the state mapping and geodesic updates are the core components that govern memory and scale preservation, utilizing this fully Softmax-free target construction maintains identical geometric properties, highlighting the massive mathematical flexibility of UGR's formulation.

We project this target distribution onto the unit hypersphere $\mathbb{S}^{K-1}$ by normalizing its Euclidean norm:
\begin{equation}
\mathbf{w}_t^{(l)} = \frac{\mathbf{e}_t^{(l)}}{\left\|\mathbf{e}_t^{(l)}\right\|_2 + \epsilon}
\label{eq:target_l2_norm}
\end{equation}
where $\epsilon = 10^{-6}$ is a numerical stability constant.

\textbf{Generalization to Alternative Similarity Kernels.} 
While we employ standard cosine similarity, $\text{sim}(x, y) = \frac{x^T y}{\|x\|_2 \|y\|_2}$, to extract bottom-up task-similarity signals, UGR's geometric formulation is fully generalizable to any bounded similarity metric or kernel. For instance, in settings where the task representations are pre-normalized, a simple inner product, $\text{sim}(x, y) = x^T y$, can be used for maximum computational efficiency and hardware friendliness. Alternatively, for highly non-linear or nested task representation manifolds, one can employ a non-linear Gaussian Radial Basis Function (RBF) kernel:
\begin{equation}
\text{sim}_{\text{RBF}}(h_t^{(l-1)}, \mu_k^{(l)}) = \exp\left(-\frac{\left\|h_t^{(l-1)} - \mu_k^{(l)}\right\|_2^2}{2\sigma^2}\right)
\end{equation}
where $\sigma > 0$ is a scale parameter. By utilizing localized non-linear similarity kernels, UGR can resolve highly overlapping or complex task boundaries while maintaining identical downstream geodesic rotation and simplex projection properties. This highlights the flexibility of the bottom-up target vector construction step, which can be tailored on-the-fly to various representation geometries.

\textbf{Quadratic Sharpening Distortion vs. Exact Born Mapping.}
We highlight a subtle but crucial design choice regarding this target construction. Under the exact square-root map (Born mapping) $s_k = \sqrt{\alpha_k}$ of Information Geometry, if we have a target distribution $\mathbf{e}_t^{(l)}$ on the probability simplex, its exact representative on the sphere should be defined as $w_{k, t}^{(l)} = \sqrt{e_{k, t}^{(l)}}$ element-wise, which naturally has unit $L_2$-norm because $\sum_k (\sqrt{e_k})^2 = \sum_k e_k = 1$. Under this exact mapping, when the state converges to the target ($\mathbf{s} \to \mathbf{w}_t^{(l)}$), the simplex weights converge linearly to the target distribution ($\alpha_k \to e_{k, t}^{(l)}$). 

In contrast, our use of $L_2$-normalization in Equation \eqref{eq:target_l2_norm} implies that when the state converges to the target, the simplex ensembling weights converge to:
\begin{equation}
\alpha_k \to (w_{k, t}^{(l)})^2 = \frac{(e_{k, t}^{(l)})^2}{\sum_j (e_{j, t}^{(l)})^2}
\end{equation}
This represents a quadratic distortion that mathematically sharpens the target distribution (amplifying dominant experts and suppressing non-dominant ones). In our empirical sandbox and real-world experiments, we find that this quadratic sharpening acts as an effective, localized task-discrimination filter that boosts classification accuracy. However, to maintain perfect mathematical alignment with the exact Fisher-Rao geodesic flow claim, one can alternatively define the target on the sphere as $w_{k, t}^{(l)} = \sqrt{e_{k, t}^{(l)}}$. We present an empirical ablation of this exact square-root target mapping, which we refer to as \texttt{UGR (Born Target)}, in Section \ref{sec:experiments}. Thus, the target construction acts as a highly customizable \textbf{``Pareto Dial of Geodesic Flow''}: standard UGR operates as a task-discriminating, boundary-sharpening accuracy-maximizer, while \texttt{UGR (Born Target)} operates as a mathematically exact, trajectory-smoothness maximizer.

\subsection{Closed-Form Geodesic Rotation Operator}
Rather than performing linear interpolation (flat EMA) which wiggles off the manifold, we perform state transitions along the shortest curved geodesic path (great-circle) connecting the previous layer state $\mathbf{s}_t^{(l-1)}$ and the bottom-up target $\mathbf{w}_t^{(l)}$. We derive a closed-form, Rodrigues-like formulation to compute this spherical interpolation (Slerp) efficiently:

\begin{enumerate}
    \item \textbf{Measure Angular Alignment:} Compute the cosine of the angle between the current state and target:
    \begin{equation}
    c_t^{(l)} = \left(\mathbf{s}_t^{(l-1)}\right)^T \mathbf{w}_t^{(l)} \in [-1, 1]
    \end{equation}
    
    \item \textbf{Handle Collinearity:} If $\left|c_t^{(l)}\right| \ge 1 - 10^{-6}$, the state and target are already aligned. To prevent division by zero, we set:
    \begin{equation}
    \mathbf{s}_t^{(l)} = \mathbf{s}_t^{(l-1)}
    \end{equation}
    
    \item \textbf{Orthonormalize Target component:} Extract the component of $\mathbf{w}_t^{(l)}$ that is orthogonal to $\mathbf{s}_t^{(l-1)}$:
    \begin{equation}
    \mathbf{v}_t^{(l)} = \mathbf{w}_t^{(l)} - c_t^{(l)} \mathbf{s}_t^{(l-1)}
    \end{equation}
    We normalize $\mathbf{v}_t^{(l)}$ to construct the orthonormal basis vector $\mathbf{u}_t^{(l)}$:
    \begin{equation}
    \mathbf{u}_t^{(l)} = \frac{\mathbf{v}_t^{(l)}}{\left\|\mathbf{v}_t^{(l)}\right\|_2}
    \end{equation}
    
    \item \textbf{Scale Angular Distance:} Compute the full geodesic angular distance $\phi_t^{(l)}$ and scale it using the step size parameter $\eta \in [0, 1]$:
    \begin{equation}
    \phi_t^{(l)} = \arccos\left(c_t^{(l)}\right)
    \end{equation}
    \begin{equation}
    \theta_t^{(l)} = \eta \cdot \phi_t^{(l)}
    \end{equation}
    
    \item \textbf{Spherical Rotation Update:} Reconstruct the updated state as a linear combination of our orthonormal basis $\{\mathbf{s}_t^{(l-1)}, \mathbf{u}_t^{(l)}\}$:
    \begin{equation}
    \mathbf{s}_t^{(l)} = \cos\left(\theta_t^{(l)}\right) \mathbf{s}_t^{(l-1)} + \sin\left(\theta_t^{(l)}\right) \mathbf{u}_t^{(l)}
    \label{eq:slerp}
    \end{equation}
\end{enumerate}

We mathematically verify that the updated state $\mathbf{s}_t^{(l)}$ is guaranteed to lie exactly on the unit hypersphere:
\begin{align}
\left\|\mathbf{s}_t^{(l)}\right\|_2^2 &= \cos^2\left(\theta_t^{(l)}\right) \left\|\mathbf{s}_t^{(l-1)}\right\|_2^2 + \sin^2\left(\theta_t^{(l)}\right) \left\|\mathbf{u}_t^{(l)}\right\|_2^2 \nonumber \\
&\quad + 2\cos\left(\theta_t^{(l)}\right)\sin\left(\theta_t^{(l)}\right) \left(\mathbf{s}_t^{(l-1)}\right)^T \mathbf{u}_t^{(l)}
\end{align}
Since $\left\|\mathbf{s}_t^{(l-1)}\right\|_2^2 = 1$, $\left\|\mathbf{u}_t^{(l)}\right\|_2^2 = 1$, and $\left(\mathbf{s}_t^{(l-1)}\right)^T \mathbf{u}_t^{(l)} = 0$ by construction, this simplifies to:
\begin{equation}
\left\|\mathbf{s}_t^{(l)}\right\|_2^2 = \cos^2\left(\theta_t^{(l)}\right) + \sin^2\left(\theta_t^{(l)}\right) = 1
\end{equation}
Thus, our closed-form update is exactly norm-preserving, eliminating representation distortions without needing costly numerical solvers.

\textbf{Computational Complexity and Scalability.} From a computational perspective, the closed-form geodesic rotation operator is exceptionally efficient. Since Slerp is composed entirely of element-wise dot-products, vector additions, and scalar trigonometric functions, it has a computational complexity of strictly $\mathcal{O}(K)$ where $K$ is the number of expert adapters. Unlike unconstrained flat-space methods that may require expensive matrix operations or high-dimensional projections, UGR scales linearly and effortlessly with the size of the expert pool, making it highly viable for large-scale multi-expert serving architectures.

We formally analyze the computational complexity and operations count of UGR compared to existing baselines in Table \ref{tab:complexity}. While SABLE, Momentum-Merge, and UGR all share the same asymptotic complexity of $\mathcal{O}(KD)$ where $D$ is the hidden representation dimension, they differ significantly in their operation constants:
\begin{itemize}
    \item \textbf{Biochemical ODE Solver Bottleneck:} ChemMerge requires virtual-time numerical integration steps (Euler or Heun's method) over $N_{\text{steps}}$ steps to update its continuous biochemical concentration states. Each step requires coordinate-wise divisions and scaling, yielding a complexity of $\mathcal{O}(KD + N_{\text{steps}} K)$ with a massive constant factor (e.g., $1 + 2 N_{\text{steps}} K$ divisions/trig ops), which creates a severe serving latency bottleneck.
    \item \textbf{Softmax-Free Geodesic Efficiency:} Our standard UGR replaces numerical ODE solvers with a single closed-form geodesic rotation, requiring only two transcendental/trig operations (arccos for torque calculation and sin/cos for Slerp). Crucially, the fully Softmax-free variant of UGR completely eliminates exponentiations by using a simple ReLU and $L_1$-normalization on similarity scores, scaling to massive expert pools with virtually zero computational overhead.
\end{itemize}

\begin{table}[h]
\centering
\caption{Computational complexity and operations count per layer update across different ensembling methods. $K$ denotes the number of experts, $D$ is the hidden representation dimension, and $N_{\text{steps}}$ is the number of virtual-time integration steps for the continuous biochemical ODE solver (ChemMerge).}
\label{tab:complexity}
\vskip 0.1in
\begin{footnotesize}
\begin{sc}
\setlength{\tabcolsep}{1.1pt}
\begin{tabular}{lccc}
\toprule
Method & Complexity & Exp. & Div / Trig \\
\midrule
SABLE & $\mathcal{O}(KD)$ & $K$ & $1$ \\
Momentum-Merge & $\mathcal{O}(KD)$ & $K$ & $1$ \\
ChemMerge (ODE) & $\mathcal{O}(KD + N_{\text{steps}} K)$ & $K$ & $1 + 2 N_{\text{steps}} K$ \\
\midrule
\textbf{UGR (Ours)} & $\mathcal{O}(KD)$ & $K$ & $2$ \\
\textbf{UGR (SF)} & $\mathcal{O}(KD)$ & $0$ & $1$ \\
\bottomrule
\end{tabular}
\end{sc}
\end{footnotesize}
\end{table}

\subsection{Torque-Driven Adaptive Agility}
A critical challenge in stateful routing is the trade-off between smoothing out representation noise (stability) and responding rapidly to task switches (plasticity). In UGR, this is resolved via \textbf{Torque-Driven Adaptive Agility}. 

In physics, torque is the rotational force that causes an object to rotate. In UGR, we define the "representational torque" as the angular distance $\phi_t^{(l)} = \arccos(c_t^{(l)})$. When the incoming activations are highly aligned with the current state (stable task stream), the alignment $c_t^{(l)} \to 1$, which makes the torque $\phi_t^{(l)} \to 0$. This causes the update step size $\theta_t^{(l)}$ to vanish, locking the ensembling weights and suppressing noise-induced jitter.

Conversely, when a sudden task switch occurs, the incoming activation signal diverges sharply from the current routing state. The alignment $c_t^{(l)} \to 0$, causing the torque $\phi_t^{(l)}$ to explode. This automatically increases the step size $\theta_t^{(l)} = \eta \phi_t^{(l)}$, allowing the state vector to rotate rapidly toward the new target expert. 

\textbf{First-Order Control-Theoretic Perspective.} From a control-theory perspective, our Torque-Driven Agility operates as a \textbf{first-order non-linear dynamical system with non-linear damping}. Specifically, the representational torque $\phi_t^{(l)}$ directly scales the \emph{angular velocity} (the step-size $\theta_t^{(l)}$) rather than the angular acceleration. This first-order formulation represents a critical control-theoretic advantage over second-order inertial or momentum-based methods (such as Momentum-Merge or ChemMerge): by scaling velocity directly with angular distance and having no second-order acceleration terms, UGR's trajectories are mathematically guaranteed to completely avoid overshoot, oscillation, or accumulation of kinetic momentum. This self-regulating property is precisely why UGR can achieve high-agility transitions under task switches while maintaining near-perfect, non-oscillating stability when the stream is stationary. This physics-inspired feedback loop operates training-free, solving representational hysteresis without hand-crafted thresholding.

\textbf{Torque Dynamics in Large Expert Pools and Concentration of Measure.} A subtle mathematical question arises when scaling UGR to extremely large expert pools (i.e., $K \gg 10^2$). In high-dimensional vector spaces, the \emph{concentration of measure} on the sphere $\mathbb{S}^{K-1}$ implies that two randomly selected unit vectors are almost always nearly orthogonal with extremely high probability. Under such conditions, if the routing state $\mathbf{s}_t^{(l-1)}$ and bottom-up target $\mathbf{w}_t^{(l)}$ behave as independent random unit vectors, their alignment cosine $c_t^{(l)} = \langle \mathbf{s}_t^{(l-1)}, \mathbf{w}_t^{(l)} \rangle$ would concentrate heavily around zero. This would cause the representational torque $\phi_t^{(l)} = \arccos(c_t^{(l)})$ to degenerate, remaining nearly constant at $\pi/2$, which would render Torque-Driven Agility insensitive to the true task boundary structure.

Fortunately, in practical Mixture-of-Experts served workloads, task representations are not randomly distributed; they reside on highly structured, low-dimensional semantic manifolds. Furthermore, we can elegantly mitigate any high-dimensional measure concentration by routing over a local, low-dimensional sub-manifold. Specifically, by extracting a top-$k$ active expert subset (where $k \ll K$, e.g., $k \in \{2, 4\}$) at each query step and restricting the geodesic rotation update strictly to the $(k-1)$-dimensional spherical sub-manifold of those active experts, the alignment cosine $c_t^{(l)}$ remains highly sensitive, and the self-regulating torque feedback loop is fully preserved. We detail this local geodesic routing extension as a scalable blueprint for massive expert served pools in Section \ref{appendix:high_dim_torque} of the Appendix.

\subsection{Spatial-Temporal Geodesic Coupling}
To maintain temporal coherence across sequential queries, we propagate the ensembling state smoothly across sample boundaries. Rather than resetting the state at the start of each query (which forces the system to struggle against artificial early-layer uniform priors), we initialize the first adapted layer's boundary condition of the current query using the final layer's state of the previous query.

For $t = 1$, we initialize the boundary state uniformly:
\begin{equation}
\mathbf{s}_1^{(L_{\text{frozen}})} = \frac{1}{\sqrt{K}} \mathbf{1} \implies \boldsymbol{\alpha}_1^{(L_{\text{frozen}})} = \frac{1}{K} \mathbf{1}
\end{equation}
For $t > 1$, we apply the coupling:
\begin{equation}
\mathbf{s}_t^{(L_{\text{frozen}})} = \mathbf{s}_{t-1}^{(L)}
\end{equation}
This creates a continuous 2D spatial-temporal geodesic trajectory across layers and sequential inputs, eliminating transition shocks.

\textbf{Mathematical and Architectural Justification of Cross-Layer Coupling.} 
One might naturally question why the state at the initial adapted layer of the current query ($\mathbf{s}_t^{(L_{\text{frozen}})}$) is initialized using the final adapted layer's state of the previous query ($\mathbf{s}_{t-1}^{(L)}$), rather than a standard layer-wise coupling scheme where each layer propagates its own independent history (i.e., $\mathbf{s}_t^{(l)} = \mathbf{s}_{t-1}^{(l)}$). We provide two fundamental justifications for this choice, which we empirically validate in Section \ref{sec:experiments}:
\begin{enumerate}
    \item \emph{Top-Down Semantic Guidance via Mature Priors:} In deep networks, early adapted layers receive noisy, highly unrefined feature activations. Consequently, early-layer routing targets ($\mathbf{w}_t^{(L_{\text{frozen}}+1)}$) are highly susceptible to local document-level noise. Conversely, as activations propagate deeper, later layers incorporate richer semantic abstractions, making the final layer's state $\mathbf{s}_{t-1}^{(L)}$ the router's most mature, confident, and denoised task belief. Propagating this final mature state to the starting boundary of the next query acts as a top-down semantic prior that stabilizes early-layer routing, shielding it from transient activation noise.
    \item \emph{Alignment with Feedforward Information Flow:} Neural networks are feedforward systems where the input of layer $l$ depends directly on the output of layer $l-1$. Under cross-layer coupling, UGR establishes a continuous 2D path where early layers inform deeper layers within the same query, and the deep semantic prior informs the start of the next query. In contrast, pure layer-wise coupling treats each layer as an independent temporal chain. This decoupling breaks the spatial dependencies across layers, leading to "spatial-temporal misalignment" where early layers and deep layers of the same query route toward conflicting experts based on their separate historical averages.
\end{enumerate}

\begin{algorithm}[t]
\caption{Unitary Geodesic Routing (UGR)}
\label{alg:ugr}
\begin{algorithmic}[1]
\REQUIRE Sequential stream of queries $\{X_t\}_{t=1}^N$, task centroids $\{\mu_k^{(l)}\}$, inertia $\eta$, temperature $\tau$.
\STATE Initialize $\mathbf{s}_1^{(L_{\text{frozen}})} = \frac{1}{\sqrt{K}} \mathbf{1}$.
\FOR{each query $X_t$ in stream}
    \STATE Pass $X_t$ through frozen layers to obtain $h_t^{(L_{\text{frozen}})}$.
    \IF{$t > 1$}
        \STATE Initialize state boundary condition: $\mathbf{s}_t^{(L_{\text{frozen}})} = \mathbf{s}_{t-1}^{(L)}$.
    \ENDIF
    \FOR{each adapted layer $l = L_{\text{frozen}}+1, \dots, L$}
        \STATE Extract hidden state $h_t^{(l-1)}$.
        \STATE Compute bottom-up similarity target $\mathbf{e}_t^{(l)}$ via Softmax with temperature $\tau$.
        \STATE Project target onto sphere: $\mathbf{w}_t^{(l)} = \mathbf{e}_t^{(l)} / (\|\mathbf{e}_t^{(l)}\|_2 + \epsilon)$.
        \STATE Compute alignment $c_t^{(l)} = \left(\mathbf{s}_t^{(l-1)}\right)^T \mathbf{w}_t^{(l)}$.
        \IF{$\left|c_t^{(l)}\right| \ge 1 - 10^{-6}$}
            \STATE $\mathbf{s}_t^{(l)} = \mathbf{s}_t^{(l-1)}$
        \ELSE
            \STATE Compute orthogonal target component: $\mathbf{v}_t^{(l)} = \mathbf{w}_t^{(l)} - c_t^{(l)} \mathbf{s}_t^{(l-1)}$.
            \STATE Normalize orthogonal component: $\mathbf{u}_t^{(l)} = \mathbf{v}_t^{(l)} / \|\mathbf{v}_t^{(l)}\|_2$.
            \STATE Calculate step size with torque: $\theta_t^{(l)} = \eta \cdot \arccos\left(c_t^{(l)}\right)$.
            \STATE Rotate state: $\mathbf{s}_t^{(l)} = \cos\left(\theta_t^{(l)}\right) \mathbf{s}_t^{(l-1)} + \sin\left(\theta_t^{(l)}\right) \mathbf{u}_t^{(l)}$.
        \ENDIF
        \STATE Compute ensembling weights: $\alpha_{k, t}^{(l)} = \left(s_{k, t}^{(l)}\right)^2$.
        \STATE Blend expert activations via Eq. \eqref{eq:forward} to obtain $h_t^{(l)}$.
    \ENDFOR
    \STATE Persist final layer state $\mathbf{s}_t^{(L)}$ to initialize next query.
\ENDFOR
\end{algorithmic}
\end{algorithm}
